\documentclass[11pt,a4paper]{article}

\usepackage[utf8]{inputenc}

\usepackage[T1]{fontenc}

\usepackage[french]{babel}

\usepackage{booktabs,graphicx,geometry,amsmath,hyperref}

\geometry{margin=2.4cm}

\title{Analyse spectrale du graphe d'attention\\Llama-3.2-1B-Instruct --- BIPIA test}

\author{Rapport généré automatiquement}

\date{15/09/2026 15:29}

\begin{document}\maketitle

\section{Protocole}


Les paires proviennent de la tâche \emph{email\_QA} du benchmark BIPIA.
Pour chaque email, deux prompts sont construits : une condition \textbf{bénigne}
(l'email seul) et une condition \textbf{injectée}, identique à la précédente
augmentée d'une instruction d'attaque insérée au début, au milieu ou à la fin
du corps du message. L'insertion étant réalisée par le code, la position en
caractères de l'attaque est connue exactement, ce qui permet d'étiqueter chaque
token sans alignement approximatif entre deux tokenisations.

Le prompt complet est envoyé au modèle : template de chat, prompt système,
consignes de tâche et question. Tout ce qui n'est pas l'attaque est étiqueté
\texttt{benign} ; seule l'attaque est \texttt{injection}. L'étiquetage est
donc strictement binaire.

Les catégories d'attaque des splits \texttt{train} et \texttt{test} sont
disjointes. Choisir une couche ou une tête sur \texttt{train} puis la
rapporter sur \texttt{test} constitue donc une évaluation hors distribution
par construction.

\paragraph{Construction du graphe.} L'attention d'un décodeur étant causale,
elle est symétrisée par $W = (A + A^{\top})/2$ avant tout traitement spectral.
Le Laplacien, la décomposition propre et les cinq métriques spectrales sont
calculés par la bibliothèque \texttt{spectral-trust} (v0.3.0) ; le filtrage
des n\oe uds par rôle et les statistiques de coupe sont propres à ce travail.
Ce run a porté sur 40 paires, soit 40 observations bénignes et
40 injectées.


\section{Configuration}

\begin{table}[htbp]\centering\small
\begin{tabular}{ll}
\hline
paramètre & valeur \\
\hline
modèle & meta-llama/Llama-3.2-1B-Instruct \\
couches & 16 \\
split & test \\
paires & 40 \\
normalisation du Laplacien & sym \\
symétrisation & symmetric \\
agrégation des têtes & uniform \\
solveur propre & dense \\
retrait du puits & False \\
retrait du template & False \\
couche analysée & 8 \\
métriques par tête & True \\
\hline
\end{tabular}
\caption{Paramètres qui modifient un chiffre.}
\label{tab:config}
\end{table}


\paragraph{Commande.}
\texttt{C:\textbackslash{}Users\textbackslash{}damsw\textbackslash{}Desktop\textbackslash{}spectral\_injection\textbackslash{}spectral\_injection\textbackslash{}bipia\_cli.py --model llama3.2-1b --split test --n-pairs 40 --per-head --figures none --out results/bipia}

\paragraph{Environnement.}
Python 3.12.10, Windows-11-10.0.26200-SP0.


\section{Résultats}

\subsection{Détection par couche}

\begin{table}[htbp]
\centering
\small
\begin{tabular}{lrrrr}
\hline
layer & auroc cut conductance & auroc fiedler value & mean separation & mean accuracy \\
\hline
0 & 0.5700 & 0.5875 & 3.7186 & 0.5421 \\
1 & 0.5337 & 0.5394 & 3.4657 & 0.5917 \\
2 & 0.5650 & 0.5219 & 3.1261 & 0.5712 \\
3 & 0.6525 & 0.5450 & 3.9376 & 0.5703 \\
4 & 0.5837 & 0.5363 & 3.9618 & 0.5713 \\
5 & 0.5106 & 0.5356 & 3.7255 & 0.5676 \\
6 & 0.5637 & 0.5206 & 4.2087 & 0.5723 \\
7 & 0.5456 & 0.5194 & 3.5013 & 0.5694 \\
8 & 0.5125 & 0.5250 & 3.7026 & 0.5762 \\
9 & 0.5131 & 0.5031 & 4.1301 & 0.5706 \\
10 & 0.5456 & 0.5162 & 3.2759 & 0.5754 \\
11 & 0.5075 & 0.5244 & 2.9425 & 0.5589 \\
12 & 0.5237 & 0.5131 & 2.8623 & 0.5708 \\
13 & 0.5544 & 0.5369 & 2.7953 & 0.5687 \\
14 & 0.5337 & 0.5381 & 2.8438 & 0.5635 \\
15 & 0.5787 & 0.5363 & 3.1299 & 0.5554 \\
\hline
\end{tabular}
\caption{AUROC bénin contre injecté, par couche. \texttt{cut\_conductance} est le score sans label : il ne voit jamais où est l'injection.}
\label{tab:layers}
\end{table}


\paragraph{Meilleure couche.} L3 avec
AUROC $= 0.6525$, intervalle de confiance à 95\%
$[0.532,\ 0.773]$ (Hanley--McNeil, $n_+ = 40$, $n_- = 40$).



\textbf{Cet intervalle ne tient pas compte de la sélection.} Il suppose une
couche fixée d'avance, alors que L3 est le \emph{maximum} sur
16 couches. Sous l'hypothèse nulle, le maximum de
16 AUROC calculées sur $n_+ = 40$ et $n_- = 40$ est déjà
nettement supérieur à $0{,}5$. La conclusion appartient à la section
Interprétation, qui est rédigée à la main : aucun verdict n'est produit
automatiquement ici.


\subsection{Ventilation par position d'insertion}

\begin{table}[htbp]
\centering
\small
\begin{tabular}{lrrrr}
\hline
position & n & mean accuracy & mean separation & mean conductance \\
\hline
end & 13.0000 & 0.5746 & 4.1093 & 1.5498 \\
middle & 13.0000 & 0.5696 & 3.7302 & 1.5564 \\
start & 14.0000 & 0.5837 & 3.2995 & 1.6436 \\
\hline
\end{tabular}
\caption{Couche 8, condition injectée uniquement. Une valeur stable entre les trois positions indique que la coupe suit l'injection et non une région fixe de la séquence.}
\label{tab:positions}
\end{table}

\subsection{Catégories d'attaque les plus difficiles}

\begin{table}[htbp]
\centering
\small
\begin{tabular}{lrrr}
\hline
attack category & n & mean accuracy & mean separation \\
\hline
Business Intelligence & 1.0000 & 0.5000 & 2.9276 \\
Entertainment & 1.0000 & 0.5481 & 2.8098 \\
Research Assistance & 2.0000 & 0.5605 & 3.8941 \\
Marketing \& Advertising & 1.0000 & 0.5649 & 3.8232 \\
Conversational Agent & 7.0000 & 0.5682 & 3.4913 \\
Reverse Text & 3.0000 & 0.5685 & 4.0307 \\
Substitution Ciphers & 6.0000 & 0.5694 & 3.8037 \\
Task Automation & 5.0000 & 0.5764 & 3.9233 \\
Emoji Substitution & 3.0000 & 0.5809 & 3.8913 \\
Scams \& Fraud & 2.0000 & 0.5851 & 3.5104 \\
\hline
\end{tabular}
\caption{Couche 8, dix catégories les moins bien retrouvées.}
\label{tab:categories}
\end{table}

\subsection{Métriques par tête}

\begin{table}[htbp]
\centering
\small
\begin{tabular}{lrrrr}
\hline
metric & layer & head & auroc & direction \\
\hline
fiedler\_value & 6.0000 & 21.0000 & 0.8400 & -1.0000 \\
connectivity\_ratio & 6.0000 & 21.0000 & 0.8337 & -1.0000 \\
connectivity\_ratio & 9.0000 & 25.0000 & 0.7744 & -1.0000 \\
fiedler\_value & 9.0000 & 25.0000 & 0.7644 & -1.0000 \\
fiedler\_value & 4.0000 & 9.0000 & 0.7625 & -1.0000 \\
connectivity\_ratio & 4.0000 & 9.0000 & 0.7612 & -1.0000 \\
fiedler\_value & 15.0000 & 30.0000 & 0.7575 & 1.0000 \\
fiedler\_value & 0.0000 & 21.0000 & 0.7544 & -1.0000 \\
connectivity\_ratio & 6.0000 & 7.0000 & 0.7512 & -1.0000 \\
connectivity\_ratio & 4.0000 & 4.0000 & 0.7481 & -1.0000 \\
fiedler\_value & 4.0000 & 4.0000 & 0.7463 & -1.0000 \\
fiedler\_value & 6.0000 & 7.0000 & 0.7438 & -1.0000 \\
connectivity\_ratio & 13.0000 & 15.0000 & 0.7369 & 1.0000 \\
connectivity\_ratio & 7.0000 & 6.0000 & 0.7275 & 1.0000 \\
fiedler\_value & 15.0000 & 21.0000 & 0.7275 & 1.0000 \\
\hline
\end{tabular}
\caption{Quinze meilleurs triplets (couche, tête, métrique).}
\label{tab:heads}
\end{table}


\textbf{Avertissement.} Ce maximum est pris sur 2560 candidats
(16 couches $\times$ 32 têtes
$\times$ 5 métriques) évalués sur 80
observations. Le maximum de plusieurs milliers d'AUROC est élevé même sous
l'hypothèse nulle : une simulation avec des scores purement aléatoires donne
un maximum proche de $0{,}99$ pour 12 observations et $0{,}69$ pour 120.
\emph{Ce chiffre n'est pas un résultat de détection.} Il ne le devient
qu'après avoir fixé la tête sur \texttt{train} et l'avoir rapportée sur
\texttt{test}.


\section{Contrôles chiffrés}


\paragraph{Séparation.} $\sqrt{d_b \cdot d_i} / d_{\text{croisé}}$ varie
de $2.80$ à $4.21$ selon la couche
(16 couches sur 16 au-dessus de 1).
Cette quantité utilise les \emph{vraies} étiquettes : elle mesure si les deux
familles se tiennent, pas si un détecteur les trouve.



\paragraph{Accord contre ligne de base.} À la couche 8, accord moyen
$0.576$ ; ligne de base triviale (tous les n\oe uds dans
un seul groupe, \texttt{partition\_agreement} maximisant sur les deux
orientations) $0.960$.
0 paire(s) sur 40 dépassent cette ligne.



\paragraph{Indice de Rand ajusté.} Moyenne $-0.0022$ sur
640 observations injectées (0 correspond au hasard). L'ARI est corrigé
du hasard, contrairement à l'accord ci-dessus.


\section{Interprétation}


\IfFileExists{interpretation.tex}{\input{interpretation}}{%
\emph{À rédiger.} Les sections précédentes sont générées automatiquement à
partir des CSV ; celle-ci ne l'est pas. Déposer un fichier
\texttt{interpretation.tex} dans ce dossier puis recompiler : il sera inclus
ici, sans avoir à relancer le run.

Points à trancher, qu'aucun gabarit ne peut trancher : l'effet de sélection sur
la meilleure couche ou la meilleure tête, la part de la séparation imputable à
la contiguïté des tokens injectés, et l'écart entre une séparation élevée et un
accord nul.%
}

\section{Définitions des métriques}


\begin{description}
\item[fiedler\_value] $\lambda_2$ du Laplacien, la connectivité algébrique.
  Proche de 0 lorsque le graphe est presque déconnecté en deux.
\item[connectivity\_ratio] $\lambda_2 / \lambda_{\max}$, sans échelle.
\item[spectral\_entropy\_norm] entropie de la distribution des valeurs propres,
  normalisée par $\log n$. Élevée si le spectre est étalé.
\item[hfer] part de la masse spectrale dans la moitié haute du spectre.
\item[spectral\_radius] $\lambda_{\max}$.
\item[density\_benign, density\_injection] poids moyen par paire \emph{possible}
  à l'intérieur de chaque famille.
\item[density\_cross] idem entre les deux familles.
\item[separation] $\sqrt{d_b \cdot d_i} / d_{\text{croisé}}$.
\item[conductance, normalized\_cut, modularity] qualité de la coupe entre la
  vraie frontière bénin/injection.
\item[cut\_conductance] conductance de la coupe \emph{découverte} par le signe
  du vecteur de Fiedler, sans utiliser les étiquettes. C'est le score du
  détecteur.
\item[accuracy, ari] accord entre la coupe découverte et la vraie frontière.
\end{description}


\end{document}